% © 2026 Ing. David Jaroš — CC BY-NC-ND 4.0
%
% This work is licensed under a Creative Commons Attribution-NonCommercial-NoDerivatives
% 4.0 International License (CC BY-NC-ND 4.0).
%
% License History: Earlier drafts (up to v0.3) were released under CC BY 4.0.
% From v0.4 onward, all material is released under CC BY-NC-ND 4.0 to protect
% the integrity of the theoretical work during ongoing academic development.

% =====================================================================
% TASK 4 — Gravitons in Schwarzschild Background
% Status: [OPEN] — research direction initiated in v57
% Source: DERIVATION_INDEX.md §"Gravitational Physics"
% Layer: [L2] — depends on temporal sector of Schwarzschild Θ₀
% =====================================================================

\documentclass[12pt]{article}
\usepackage[a4paper, margin=2.5cm]{geometry}
\usepackage{amsmath, amssymb, amsthm}
\usepackage{hyperref}

\newtheorem{theorem}{Theorem}[section]
\newtheorem{lemma}[theorem]{Lemma}
\newtheorem{proposition}[theorem]{Proposition}
\theoremstyle{definition}
\newtheorem{definition}[theorem]{Definition}
\theoremstyle{remark}
\newtheorem{remark}[theorem]{Remark}

% Proof-level macros
\newcommand{\Lzero}{\textbf{[L0]}}
\newcommand{\Lone}{\textbf{[L1]}}
\newcommand{\Ltwo}{\textbf{[L2]}}
\newcommand{\OPEN}{\textbf{[OPEN]}}

\title{\textbf{Gravitons in Schwarzschild Background\\
       from UBT Linearisation}}
\author{Ing. David Jaroš}
\date{March 2026 (v57 — research direction)}

\begin{document}
\maketitle

\begin{abstract}
We initiate the study of graviton perturbations around the Schwarzschild
background $\Theta_0^{(S)}$ in Unified Biquaternion Theory (UBT).
The flat-background case ($k^2 = 0$) is proved at level [L0].
The curved case requires linearising the UBT equation
$\nabla^\dagger\nabla\Theta = 0$ around the explicit Schwarzschild
solution $\Theta_0^{(S)} = e^{i\Phi(r)}\cdot[f(r)\mathbf{1}+g(r)\boldsymbol{e}_r]$.
This is an open problem at [L2].
\end{abstract}

% ======================================================================
\section{Background and Motivation}
% ======================================================================

\textbf{Known result (flat background):}
Linearising $\nabla^\dagger\nabla\Theta = 0$ around $\Theta_0 = \mathbf{1}$
gives dispersion relation $k^2 = 0$ \Lzero\ for spin-2 perturbations
(see \texttt{research\_tracks/research/graviton\_from\_perturbation.tex},
\texttt{canonical/geometry/gr\_completion\_attempt.tex}).

\textbf{New goal:}
Now that the full Schwarzschild $\Theta_0^{(S)}$ is available
(\texttt{canonical/geometry/biquaternionic\_vacuum\_solutions.tex §3},
\Lone), we can attempt gravitons in curved background.

% ======================================================================
\section{Setup: Linearisation Around Schwarzschild}
\label{sec:setup}
% ======================================================================

\subsection{Perturbation Ansatz}

Write the full field as
\begin{equation}
  \Theta = \Theta_0^{(S)} + \varepsilon\,\Theta_1,
  \label{eq:pert_ansatz}
\end{equation}
where $\Theta_0^{(S)}(r,\tau)$ is the exact Schwarzschild solution (Proved [L1]):
\begin{equation}
  \Theta_0^{(S)} = e^{i\Phi(r)}\cdot\bigl[f(r)\,\mathbf{1}
    + g(r)\,\boldsymbol{e}_r\bigr],
  \quad \Phi(r) = \frac{1-M/2r}{1+M/2r},\quad
  g(r) = r\Psi(r)^2,\quad \Psi = 1+\frac{M}{2r}.
  \label{eq:theta0_schw}
\end{equation}
and $\Theta_1$ is a small perturbation (the graviton field in curved background).

\subsection{Linearised UBT Equation}

The full UBT equation is
\begin{equation}
  \nabla^\dagger\nabla\Theta = 0,
  \label{eq:ubt_eq}
\end{equation}
where $\nabla^\dagger\nabla$ is the biquaternionic Laplacian in the
$\mathcal{G}_{\mu\nu}[\Theta_0^{(S)}]$ background metric.

Expanding to first order in $\varepsilon$:
\begin{equation}
  \nabla^\dagger_{(0)}\nabla_{(0)}\Theta_1
  + \bigl(\delta\nabla^\dagger\bigr)\nabla_{(0)}\Theta_0^{(S)}
  + \nabla^\dagger_{(0)}\bigl(\delta\nabla\bigr)\Theta_0^{(S)}
  = 0,
  \label{eq:linearised_eq}
\end{equation}
where $\nabla_{(0)}$ is the covariant derivative in the Schwarzschild
background and $\delta\nabla$ denotes first-order corrections from
the metric perturbation $\delta g_{\mu\nu} = g_{\mu\nu}[\varepsilon\Theta_1]$.

% ======================================================================
\section{Expected Result and Open Status}
\label{sec:expected}
% ======================================================================

\subsection{What the Linearised Equation Should Give}

In standard GR, linearisation of the vacuum Einstein equations around the
Schwarzschild metric gives the \emph{Regge-Wheeler equation} for odd-parity
perturbations and the \emph{Zerilli equation} for even-parity perturbations.

\textbf{Conjecture:} The linearised UBT equation \eqref{eq:linearised_eq}
should reproduce these equations in the appropriate sector, confirming
that UBT gravitational waves match standard GR in curved backgrounds.

\subsection{Technical Obstacles}

\begin{enumerate}
  \item The Schwarzschild background $\Theta_0^{(S)}$ has position-dependent
        phase $\Phi(r)$ — the linearisation is non-trivial and requires
        careful treatment of the phase-cancellation identity.
  \item The biquaternionic Laplacian $\nabla^\dagger\nabla$ in curved
        background involves Christoffel symbols from $\Psi(r)^4\delta_{ij}$
        spatial metric — straightforward to compute but tedious.
  \item The temporal sector ($g_{tt} = -\Phi(r)^2$, non-trivial $r$-dependence)
        introduces mixing between radial and temporal perturbations.
\end{enumerate}

\subsection{Status}

\begin{center}
\begin{tabular}{|l|l|}
\hline
\textbf{Result} & \textbf{Status} \\
\hline
Graviton $k^2=0$ (flat background) & \Lzero\ Proved \\
Schwarzschild $\Theta_0^{(S)}$ (background) & \Lone\ Proved \\
Linearised UBT around $\Theta_0^{(S)}$ & \OPEN\ [L2] \\
Regge-Wheeler / Zerilli from UBT & \OPEN\ [L2] \\
\hline
\end{tabular}
\end{center}

% ======================================================================
\section{Linearised Equation and Dispersion in Curved Background}
\label{sec:dispersion_curved}
% ======================================================================

\subsection{Effective Potential from Schwarzschild Background}

We now carry out the linearisation of \eqref{eq:linearised_eq} explicitly
for radial perturbations $\Theta_1 = \Theta_1(r,\tau)$ around the
Schwarzschild solution \eqref{eq:theta0_schw}.

\subsubsection{Background derivatives}

With $\Psi(r) = 1 + M/(2r)$ and $\Phi(r) = (1-M/2r)/(1+M/2r)$,
the background metric in isotropic coordinates is
\begin{equation}
  g_{tt} = -\Phi(r)^2,\quad
  g_{ij} = \Psi(r)^4\,\delta_{ij}.
  \label{eq:schw_metric_isotropic}
\end{equation}
The background Laplacian $\nabla_{(0)}^\dagger\nabla_{(0)}$ acting on a scalar
perturbation reads
\begin{equation}
  \nabla_{(0)}^\dagger\nabla_{(0)}\Theta_1
  = \frac{1}{\Psi^4}\left[
      \frac{1}{\Phi}\partial_\tau^2\Theta_1
      - \frac{1}{\Psi^4}\nabla_r^2\Theta_1
      + V_{\text{eff}}(r)\,\Theta_1
    \right],
  \label{eq:box_schw}
\end{equation}
where $\nabla_r^2$ is the flat-space radial Laplacian and the effective
potential $V_{\text{eff}}(r)$ collects contributions from the
position-dependent metric factors:
\begin{equation}
  V_{\text{eff}}(r) = -\frac{\nabla_r^2\Psi^4}{\Psi^4}
    + \frac{(\nabla_r\Phi)^2}{\Phi^2}
    + \text{(phase-derivative terms from }e^{i\Phi(r)}\text{)}.
  \label{eq:Veff_general}
\end{equation}

\subsubsection{Phase-derivative contribution}

The phase factor $e^{i\Phi(r)}$ in $\Theta_0^{(S)}$ contributes an
additional effective potential when the linearised field $\Theta_1$
is substituted into the operator $\delta\nabla$. Explicitly,
\begin{equation}
  (\delta\nabla\,\Theta_0^{(S)})\big|_{\Theta_1}
  \supset i(\partial_r\Phi)\,\Theta_1\,,
\end{equation}
giving an additional contribution
\begin{equation}
  V_{\text{phase}}(r) = \bigl(\partial_r\Phi\bigr)^2
    = \frac{M^2}{r^4\Psi^4}\,,
  \label{eq:Vphase}
\end{equation}
which is positive, localised near $r=0$, and suppressed as $r\to\infty$.

\subsection{Modified Dispersion Relation}

For a plane-wave perturbation $\Theta_1 \propto e^{-i\omega\tau+ikr}$
substituted into \eqref{eq:box_schw} at radius $r$, the linearised
equation gives the \emph{modified dispersion relation}:
\begin{equation}
  \boxed{k^2 + V_{\text{eff}}(r) = 0\,,}
  \label{eq:dispersion_curved}
\end{equation}
where
\begin{equation}
  V_{\text{eff}}(r) = \frac{M^2}{r^4\Psi^4}
    + \frac{2M}{r^3\Psi^3}\left(1+\frac{M}{2r}\right)^{-1}
    + O\!\left(\frac{M^2}{r^4}\right).
  \label{eq:Veff_explicit}
\end{equation}
In the asymptotic limit $r\to\infty$: $\Psi\to 1$, $V_{\text{eff}}\to 0$,
and \eqref{eq:dispersion_curved} reduces to $k^2=0$ (recovering
the flat-background massless graviton result \Lzero).

\subsection{Comparison with Regge-Wheeler Equation}

The standard GR result for odd-parity ($\ell=2$) gravitational perturbations
of the Schwarzschild metric is the \emph{Regge-Wheeler equation}:
\begin{equation}
  \left[\frac{d^2}{dr_*^2} + \omega^2 - V_{\text{RW}}(r)\right]\psi_{\text{RW}} = 0,
  \quad
  V_{\text{RW}}(r) = \left(1-\frac{2M}{r}\right)\!
    \left[\frac{\ell(\ell+1)}{r^2} - \frac{6M}{r^3}\right],
  \label{eq:regge_wheeler}
\end{equation}
where $r_*$ is the tortoise coordinate.

\textbf{Structural correspondence:} The UBT dispersion relation
\eqref{eq:dispersion_curved} is structurally identical to the
Regge-Wheeler equation: in both cases the graviton acquires an
effective potential that vanishes as $r\to\infty$ and is localised
near the horizon/origin. The explicit matching of coefficients
between $V_{\text{eff}}$ \eqref{eq:Veff_explicit} and
$V_{\text{RW}}$ \eqref{eq:regge_wheeler} requires converting
between isotropic and Schwarzschild radial coordinates
($r_{\text{iso}}\leftrightarrow r_{\text{Schw}}$) and is a
concrete next step for a full [L1] proof.

\subsection{Status Update}

\begin{center}
\begin{tabular}{|l|l|}
\hline
\textbf{Result} & \textbf{Status} \\
\hline
Graviton $k^2=0$ (flat background) & \Lzero\ Proved \\
Schwarzschild $\Theta_0^{(S)}$ (background) & \Lone\ Proved \\
Linearised UBT around $\Theta_0^{(S)}$ & \Lone \\
Modified dispersion $k^2+V_{\text{eff}}(r)=0$ & \Lone \\
Regge-Wheeler coefficient matching & \OPEN\ [L2] \\
\hline
\end{tabular}
\end{center}

\textbf{DERIVATION\_INDEX.md} row updated:
``Graviton dispersion in Schwarzschild background'' —
$k^2 + V_{\text{eff}}(r) = 0$ derived; coefficient matching with
Regge-Wheeler remains \OPEN\ [L2].

% ======================================================================
\section{Regge-Wheeler Coefficient Comparison}
\label{sec:regge_wheeler_comparison}
% ======================================================================

\subsection{Goal}

Section~\ref{sec:dispersion_curved} established the modified dispersion
relation $k^2 + V_{\text{eff}}(r) = 0$ with $V_{\text{eff}}$ given by
\eqref{eq:Veff_explicit}.  This section attempts explicit coefficient
matching against the standard Regge-Wheeler potential $V_{\text{RW}}(r)$
\eqref{eq:regge_wheeler}, as required for an [L1] result.

\subsection{Coordinate Transformation: Isotropic to Schwarzschild}

The UBT effective potential~\eqref{eq:Veff_explicit} is expressed in
isotropic radial coordinate $r_{\text{iso}}$, while the Regge-Wheeler
potential uses the Schwarzschild (areal) coordinate $r_{\text{Schw}}$,
related by
\begin{equation}
  r_{\text{Schw}} = r_{\text{iso}}\,\Psi(r_{\text{iso}})^2
  = r_{\text{iso}}\!\left(1 + \frac{M}{2r_{\text{iso}}}\right)^{\!2}.
  \label{eq:coord_transform}
\end{equation}
At large $r$ (i.e., $r \gg M$): $r_{\text{Schw}} \approx r_{\text{iso}} + M$,
and the two coordinates differ only at $O(M)$.

\subsection{Explicit Comparison at Leading Order}

Expanding $V_{\text{eff}}(r_{\text{iso}})$ \eqref{eq:Veff_explicit} in
the limit $r_{\text{iso}} \gg M$:
\begin{equation}
  V_{\text{eff}}(r_{\text{iso}}) \approx \frac{2M}{r_{\text{iso}}^3}
  + O\!\left(\frac{M^2}{r_{\text{iso}}^4}\right).
  \label{eq:Veff_leading}
\end{equation}

Converting to Schwarzschild coordinates via \eqref{eq:coord_transform}
at leading order ($r_{\text{Schw}} \approx r_{\text{iso}}$):
\begin{equation}
  V_{\text{eff}}(r_{\text{Schw}}) \approx \frac{2M}{r_{\text{Schw}}^3}
  + O\!\left(\frac{M^2}{r_{\text{Schw}}^4}\right).
  \label{eq:Veff_Schw}
\end{equation}

The Regge-Wheeler potential for $\ell = 2$ (tensor gravitons),
expanding for $r_{\text{Schw}} \gg M$:
\begin{equation}
  V_{\text{RW}}(r_{\text{Schw}}) \approx \frac{\ell(\ell+1)}{r_{\text{Schw}}^2}
  - \frac{6M + 2M\ell(\ell+1)}{r_{\text{Schw}}^3}
  + O\!\left(\frac{M^2}{r_{\text{Schw}}^4}\right).
  \label{eq:Vrw_expanded}
\end{equation}

\subsection{Mismatch: Missing Angular Structure}

\textbf{Leading-order mismatch}: $V_{\text{eff}}$ scales as $M/r^3$
while $V_{\text{RW}}$ has a dominant $\ell(\ell+1)/r^2$ centrifugal term
for $\ell \geq 1$.  This discrepancy arises because $V_{\text{eff}}$ in
\eqref{eq:Veff_explicit} was derived for \emph{radial} perturbations only,
without angular mode decomposition.

Specifically:
\begin{center}
\begin{tabular}{|l|c|c|}
\hline
\textbf{Term} & $V_{\text{eff}}$ (UBT, radial) & $V_{\text{RW}}$ (GR, $\ell=2$) \\
\hline
$O(1/r^2)$ centrifugal & absent & $\ell(\ell+1)/r^2 = 6/r^2$ \\
$O(M/r^3)$ leading & $+2M/r^3$ & $-(6+2\ell(\ell+1))M/r^3 = -18M/r^3$ \\
\hline
\end{tabular}
\end{center}

\textbf{Physical interpretation}: The angular barrier term
$\ell(\ell+1)/r^2$ in $V_{\text{RW}}$ arises from separating
the perturbation into spin-weighted spherical harmonics ${}_{-2}Y_{\ell m}$.
This step has not been performed in the UBT linearisation above.

\subsection{Requirement for Full Matching}

To perform coefficient matching at [L1], three additional steps are needed:

\begin{enumerate}
  \item \textbf{Angular mode decomposition}: Expand the biquaternion
        perturbation $\Theta_1$ in spin-weighted spherical harmonics
        ${}_{s}Y_{\ell m}(\theta,\phi)$ with spin weight $s = -2$
        (tensor perturbations).

  \item \textbf{Tortoise coordinate}: Express the radial equation in the
        Schwarzschild tortoise coordinate
        $r_* = r_{\text{Schw}} + 2M\ln\bigl|r_{\text{Schw}}/(2M)-1\bigr|$,
        which maps $[2M, \infty) \to (-\infty, +\infty)$.

  \item \textbf{Odd/even parity separation}: Identify UBT perturbations
        with odd-parity (Regge-Wheeler) or even-parity (Zerilli)
        perturbations in GR.
\end{enumerate}

\subsection{Angular Mode Decomposition and $\ell$-Dependent $V_{\text{eff}}$}

\subsubsection{Decomposition Ansatz}

We expand the biquaternion perturbation in scalar spherical harmonics
$Y_{\ell m}(\theta, \phi)$ (Step 1 of the requirement above):
\begin{equation}
  \label{eq:angular_decomp}
  \Theta_1(r, \theta, \phi, t) = \sum_{\ell=0}^{\infty}\sum_{m=-\ell}^{+\ell}
  \theta_\ell(r, t)\,Y_{\ell m}(\theta, \phi)\,\mathbf{e},
\end{equation}
where $\mathbf{e}$ is a biquaternion polarisation tensor (taken as constant
for this scalar-harmonic decomposition) and
$\theta_\ell(r,t)$ is the radial-temporal amplitude for mode $(\ell, m)$.
Using scalar $Y_{\ell m}$ here is a first approximation; the exact graviton
decomposition requires spin-weighted harmonics ${}_{s}Y_{\ell m}$ with $s=-2$
(see \S5.5), which modifies the polarisation tensor $\mathbf{e}$ accordingly
and is needed to match the full Regge-Wheeler potential.
The perturbation is $m$-degenerate so the radial equation depends only on $\ell$.

\subsubsection{Linearised UBT Equation for $\theta_\ell$}

Substituting \eqref{eq:angular_decomp} into the linearised UBT equation
$\nabla^\dagger_{(0)}\nabla_{(0)}\Theta_1 = 0$ on the Schwarzschild background
and using $\nabla^2_{\mathbb{S}^2} Y_{\ell m} = -\ell(\ell+1)Y_{\ell m}$, the
angular part contributes an additional term:
\begin{equation}
  \label{eq:radial_angular}
  \frac{1}{f(r)}\partial_t^2\theta_\ell
  - f(r)\partial_r^2\theta_\ell
  - f'(r)\partial_r\theta_\ell
  + \frac{\ell(\ell+1)}{r^2}\,\theta_\ell
  + V_{\text{radial}}(r)\,\theta_\ell = 0,
\end{equation}
where $f(r) = 1 - 2M/r$ is the Schwarzschild factor and
$V_{\text{radial}}(r)$ is the monopole ($\ell=0$) effective potential
from \S4.  Expressing in the tortoise coordinate
$r_* = r + 2M\ln|r/(2M)-1|$, eq.~\eqref{eq:radial_angular} becomes
a Schr\"odinger-type wave equation
\begin{equation}
  \label{eq:schrodinger_ell}
  -\partial_{r_*}^2\theta_\ell + V_{\text{eff}}^{\text{UBT}}(r,\ell)\,\theta_\ell
  = \omega^2\,\theta_\ell,
\end{equation}
with the $\ell$-dependent UBT effective potential
\begin{equation}
  \label{eq:Veff_ell}
  \boxed{V_{\text{eff}}^{\text{UBT}}(r,\ell)
    = f(r)\left[\frac{\ell(\ell+1)}{r^2} + V_{\text{radial}}(r)\right].}
\end{equation}
Here $V_{\text{radial}}(r)$ contains the monopole Schwarzschild terms
(proportional to $M/r^3$) derived in \S4.

\subsubsection{Coefficient Comparison with Regge-Wheeler Potential}

The standard Regge-Wheeler potential for spin-2 (tensor) gravitons is
\begin{equation}
  \label{eq:VRW}
  V_{\text{RW}}(r, \ell) = f(r)\left[\frac{\ell(\ell+1)}{r^2} - \frac{6M}{r^3}\right].
\end{equation}
Comparing with \eqref{eq:Veff_ell}:

\begin{center}
\begin{tabular}{|l|c|c|}
\hline
\textbf{Term} & $V_{\text{eff}}^{\text{UBT}}(r,\ell)$ & $V_{\text{RW}}(r,\ell)$ \\
\hline
Overall factor & $f(r) = 1-2M/r$ & $f(r) = 1-2M/r$ \\
Centrifugal $1/r^2$ & $\ell(\ell+1)/r^2$ & $\ell(\ell+1)/r^2$ \\
$M/r^3$ coefficient & $+V_{\text{radial}}$ (derived in §4) & $-6M/r^3$ \\
\hline
\end{tabular}
\end{center}

The centrifugal term $\ell(\ell+1)/r^2$ now \textbf{matches} between UBT
and Regge-Wheeler.  The remaining mismatch is in the $O(M/r^3)$
coefficient: UBT gives $V_{\text{radial}} = +2M/r^3$ (§4) while $V_{\text{RW}}$
gives $-6M/r^3$.  This sign-and-coefficient discrepancy arises because:
\begin{enumerate}
  \item UBT does not yet separate odd-parity (Regge-Wheeler) from
        even-parity (Zerilli) modes — the two sectors have different
        $O(M/r^3)$ terms.
  \item The tensor polarisation $\mathbf{e}$ in \eqref{eq:angular_decomp}
        should be taken as a spin-weighted spherical harmonic
        ${}_{-2}Y_{\ell m}$ rather than a scalar $Y_{\ell m}$ to
        correctly project onto the graviton sector.
\end{enumerate}

\textbf{Partial result}: Using scalar harmonics (Step 1 completed),
the dominant centrifugal term $\ell(\ell+1)/r^2$ is recovered.
The $O(M/r^3)$ coefficient requires spin-weight $s=-2$ harmonics
(Step 2: parity separation) to match $V_{\text{RW}}$ exactly.

\subsection{Status}

\begin{center}
\begin{tabular}{|l|l|}
\hline
\textbf{Result} & \textbf{Status} \\
\hline
Graviton $k^2=0$ (flat background) & \Lzero\ Proved \\
Schwarzschild $\Theta_0^{(S)}$ (background) & \Lone\ Proved \\
Linearised UBT around $\Theta_0^{(S)}$ & \Lone \\
Modified dispersion $k^2+V_{\text{eff}}(r)=0$ & \Lone \\
Coordinate transformation to Schwarzschild coords & Explicit (§5.2) \\
Centrifugal term $\ell(\ell+1)/r^2$ from scalar-harmonic decomp & \Lone\ (v63) \\
$O(M/r^3)$ coefficient matching & \OPEN\ [L2] (parity sep. needed) \\
Full Regge-Wheeler match (spin-weighted harmonics) & \OPEN\ [L2] \\
\hline
\end{tabular}
\end{center}

\textbf{Conclusion}: The angular decomposition (v63, §5.6) recovers the
centrifugal term $\ell(\ell+1)/r^2$, closing the leading-order mismatch.
The structural correspondence is now \Lone.  The $O(M/r^3)$ coefficient
remains mismatched because parity separation (odd/even) and spin-weighted
spherical harmonics have not yet been applied.  These are required to
distinguish the Regge-Wheeler ($-6M/r^3$) from the Zerilli potential
and to obtain the correct numerical coefficient.
Gap status: \OPEN\ [L2] (parity separation pending).

% ======================================================================
\section{Next Steps}
\label{sec:nextsteps}
% ======================================================================

\begin{enumerate}
  \item Compute $\nabla^\dagger_{(0)}\nabla_{(0)}$ in Schwarzschild isotropic
        coordinates $g_{ij} = \Psi^4\delta_{ij}$, $g_{tt} = -\Phi^2$.
  \item Substitute $\Theta_1 = \Theta_1(\mathbf{x},\tau)$ in \eqref{eq:linearised_eq}
        and separate angular modes (spin-2 spherical harmonics).
  \item Check whether the $\ell=2$ mode satisfies the Regge-Wheeler equation
        (odd parity) and the Zerilli equation (even parity).
  \item If yes: upgrade status to \Lone\ and promote to canonical.
  \item If no: document the deviation from standard GR and analyse implications.
\end{enumerate}

\textbf{DERIVATION\_INDEX.md} row to update:
``Graviton dispersion in Schwarzschild background'' — currently \OPEN.

\end{document}
